% !TeX root = ../../Book5.tex
% 纲要标题：### A.2 对称函数、Young 表与组合恒等式
% 纲要位置：工作记录/计划纲要.md，第 4547 行。
% #### A.2.1 半标准表与 Schur 多项式
% #### A.2.2 行插入与成对标准表
% #### A.2.3 维数恒等式与格路径计数
\section{对称函数、Young 表与组合恒等式}
\label{b5:sec:A02}

第6章用 Young 表构造 Specht 模，第7章用对称函数计算它们的特征标。这里补充表格本身的计数解释，并证明置换与成对标准表之间的对应。行表类是按行内置换得到的等价类，列交错元则属于以行表类为基的模；以下的插入算法始终操作具体的表，不对等价类作插入。

\subsection{半标准表与 Schur 多项式}
\label{b5:subsec:A02:01}

\begin{definition}\label{b5:def:semistandard-tableau}
设 \(\lambda\) 为分拆，\(N\ge1\) 为整数，\(T\) 为在 Young 图 \([\lambda]\) 各格填入 \(1,\ldots,N\) 得到的表。若每行从左到右弱递增、每列从上到下严格递增，则称 \(T\) 为\term[banbiaozhun-Young-biao]{半标准 Young 表}[semistandard Young tableau]。若数字 \(j\) 出现 \(a_j\) 次，称 \((a_1,\ldots,a_N)\) 为其重量，并对形式变量 \(x_1,\ldots,x_N\) 记 \(x^T=\prod_{j=1}^Nx_j^{a_j}\)。
\end{definition}
\symidx{\(x^T\)：半标准 Young 表的权重单项式}

\begin{proposition}\label{b5:prop:schur-tableau-expansion}
设 \(\lambda\) 为分拆，\(N\ge1\) 为整数，\(x_1,\ldots,x_N\) 为形式变量，\(s_\lambda\in\mathbb Q[x_1,\ldots,x_N]\) 为 Schur 多项式。则
\[
s_\lambda(x_1,\ldots,x_N)=\sum_T x^T,
\]
其中 \(T\) 遍历用 \(1,\ldots,N\) 填写的形状 \(\lambda\) 的半标准表，\(x^T=\prod_{j=1}^Nx_j^{a_j(T)}\)，\(a_j(T)\) 为数字 \(j\) 在 \(T\) 中出现的次数。
\end{proposition}
\begin{proof}
对半标准表 \(T\)，数字不超过 \(j\) 的格子组成一个 Young 图 \(\lambda^{(j)}\)：行弱递增保证每行是前缀，列严格递增保证上面的格子先出现。同一个数字不能出现在同列，所以
\(\lambda^{(j)}/\lambda^{(j-1)}\) 为横条。反过来，由一条
\[
\varnothing=\lambda^{(0)}\subseteq\lambda^{(1)}
\subseteq\cdots\subseteq\lambda^{(N)}=\lambda
\]
且每次增加横条的链，把第 \(j\) 步新增格子填为 \(j\)，就得到半标准表。这是互逆对应。

由\cref{b5:lem:schur-cauchy}，对另一组变量 \(z\) 有
\[
\prod_{i,j}(1-z_ix_j)^{-1}
=\sum_\lambda s_\lambda(z)s_\lambda(x).
\]
左端也等于
\(\prod_{j=1}^N\sum_{r\ge0}h_r(z)x_j^r\)。
使用\cref{b5:thm:pieri}逐次展开各 \(h_r(z)\) 的乘积，\(s_\lambda(z)\) 的系数恰为上述横条链的重量和。Schur 函数在每个次数构成一组基，比较系数就得到公式。
\end{proof}

例如形状 \((2,1)\) 使用 \(1,2\) 时只有两张半标准表，重量分别为 \((2,1)\)、\((1,2)\)，所以
\(s_{(2,1)}(x_1,x_2)=x_1^2x_2+x_1x_2^2\)。当允许 \(N\) 个数字时，代入 \(x_i=1\) 就计数全部半标准表；这一数目也等于第7章中对应 Schur 模的维数。

\subsection{行插入与成对标准表}
\label{b5:subsec:A02:02}

对数字互异、行列均严格递增的表，插入一个未出现的数字 \(a\) 的规则如下。先在第一行找到第一个严格大于 \(a\) 的数字 \(b\)，用 \(a\) 替换它，并把 \(b\) 插入下一行；若本行没有大于待插数字的项，就把该数字附到本行末尾并停止。若下一行原不存在，就新建一行。

\begin{lemma}\label{b5:lem:row-insertion}
设 \(T\) 为 Young 图上的填数表，所有数字互异且行列均严格递增，\(a\) 为未出现在 \(T\) 中的数字。把 \(a\) 插入第一行：每行用待插数字替换该行第一个严格大于它的数字，并将被替换的数字继续插入下一行；若没有更大的数字，就在该行末尾追加并停止，必要时新建一行。这个行插入过程保持行列严格递增，并且只在一个外角增加一格。给定插入后的表及该新增外角，可以唯一恢复原表和最初插入的数字。
\end{lemma}
\begin{proof}
每行的替换位置取第一个大于待插数字的位置，故替换后该行仍严格递增。设上一行在第 \(c\) 列弹出了数字 \(b\)。下一行原第 \(c\) 列若存在，其数字大于 \(b\)，所以这次替换或追加的位置 \(c'\le c\)。在第 \(c'\) 列上方，若 \(c'<c\)，数字小于原上一行第 \(c\) 列的 \(b\)；若 \(c'=c\)，它已被更小的数字替换。因此新放入的 \(b\) 大于其上方数字。它又小于被替换的旧数字，因而仍小于下方原有数字。追加时没有下方格子：上一行长度至少达到追加位置，而本行原先比该位置短，下方各行更短。故新图仍为 Young 图且新增格为外角。

逆过程从指定外角取出数字 \(b\)，删掉该格；在上一行找到最右边严格小于 \(b\) 的数字 \(a\)，用 \(b\) 替换它，再将 \(a\) 向上一行继续传递，直至离开第一行。所需位置存在，因为当前列的上方数字小于正在向上移动的数字。正向过程中，被插数字左边的数都小于它，右边的数都大于被弹数字；所以逆向选择的“最右边小于”恰找回原替换位置。逐行反转即恢复每一步，证明存在性及唯一性。
\end{proof}

下面的 Robinson--Schensted 对应\footnote{鲁宾逊（Gilbert de Beauregard Robinson，1906-1992），加拿大数学家，研究对称群表示和组合学；申斯特德（Craige Schensted，1927-2021），美国数学家、物理学家，提出行插入算法。}把置换配上一对表。
\begin{theorem}[Robinson--Schensted 对应]\label{b5:thm:rs-correspondence}
设 \(n\) 为非负整数；约定 \(S_0=\{1\}\)，空分拆的 Young 图有唯一的空标准表。置换 \(w\in S_n\) 与一对同形状、各用 \(1,\ldots,n\) 填数的标准 Young 表 \((P,Q)\) 一一对应：按 \(w_i=w(i)\) 的次序行插入，得到插入表 \(P\)，并在每步的新增格记录步数，得到记录表 \(Q\)。当 \(n=0\) 时，唯一的置换对应空表对 \((\varnothing,\varnothing)\)。
\end{theorem}
\begin{proof}
当 \(n=0\) 时，两边都只有一个对象。以下设 \(n\ge1\)。依次把 \(w_1,\ldots,w_n\) 行插入到空表，得到 \(P\)。第 \(i\) 步新增一格，在第二张图的相同位置写上 \(i\)，得到记录表 \(Q\)。前一引理保证 \(P\) 为标准表；各步形状都是 Young 图，格子右边或下边的格子必后加入，所以 \(Q\) 也是标准表。

反过来，从 \(Q\) 中找到最大数字 \(n\) 所在的外角，在 \(P\) 中从同一外角作逆插入，得到数字 \(w_n\)，并删去 \(Q\) 的这格。递减重复，得到 \(w_{n-1},\ldots,w_1\)。每个数字恰从原来填有 \(1,\ldots,n\) 的 \(P\) 中移出一次，所以输出是一个置换。引理的逐步互逆性说明两个整体过程互逆。
\end{proof}

\Cref{b5:alg:robinson-schensted}把这两个互逆过程写成可逐格执行的步骤。正向新增的格由记录表保存，逆向按记录数字从大到小删除外角。

\begin{algorithm}[htbp]
\caption{Robinson--Schensted 行插入与逆插入}
\label{b5:alg:robinson-schensted}
\begin{algorithmsteps}{置换 \(w\in S_n\)，或一对同形状、各填 \(1,\ldots,n\) 的标准 Young 表 \((P,Q)\)。}{输入为置换时输出对应的表对 \((P,Q)\)；输入为表对时输出置换 \(w\)。}
\If{输入为置换 \(w\)}
\State 置 \(P,Q\) 为空表。
\For{\(i=1,\ldots,n\)}
\State 从首行起把 \(w_i\) 插入 \(P\)：在每行替换第一个更大的数，把被替换数送往下一行；若没有更大的数，就在该行末添格并停止。记新增格为 \(c_i\)。
\State 在 \(Q\) 的同一格 \(c_i\) 填入 \(i\)。
\EndFor
\State 输出 \((P,Q)\)。
\Else
\For{\(i=n,n-1,\ldots,1\)}
\State 找到 \(Q\) 中数字 \(i\) 所在的外角，记其行为 \(r\)，并删去该格；从 \(P\) 的同一格取出数字 \(b\) 并删格。
\State 依次经过第 \(r-1,\ldots,1\) 行：找本行最右边小于 \(b\) 的数 \(a\)，用 \(b\) 替换它，再令 \(b\leftarrow a\)。最后置 \(w_i=b\)。
\EndFor
\State 输出 \(w\)。
\EndIf
\end{algorithmsteps}
\end{algorithm}

每次正向插入至多新建一行，逆向每步都向上一行移动，因此两种操作均在有限步终止；其互逆性由前述引理保证。

例如 \(w=312\) 的插入过程先得到单格 \(3\)，再把 \(1\) 插入并弹出 \(3\)，最后将 \(2\) 附到第一行。于是
\[
P=\begin{array}{cc}1&2\\3&\end{array},
\qquad
Q=\begin{array}{cc}1&3\\2&\end{array}.
\]
从 \(Q\) 的格子 \(3,2,1\) 依次逆插入，正好恢复 \(2,1,3\) 的逆读取次序。

\subsection{维数恒等式与格路径计数}
\label{b5:subsec:A02:03}

\begin{corollary}\label{b5:cor:tableau-squares}
设 \(n\) 为非负整数，对每个分拆 \(\lambda\vdash n\)，记 \(f^\lambda\) 为形状 \(\lambda\) 的标准表数，则
\[
\sum_{\lambda\vdash n}(f^\lambda)^2=n!.
\]
\end{corollary}
\begin{proof}
对固定形状，可独立选择 \(P,Q\)，共有 \((f^\lambda)^2\) 对。上一对应把所有这些对恰好与 \(n!\) 个置换对应，求和即得。
\end{proof}

复群代数分解也给出这一维数恒等式；插入法则直接给出组合证明。

\ProseParagraphBreak
两行形状 \((a,b)\)、\(a\ge b\ge1\) 还有直接的路径计数。按 \(1,\ldots,a+b\) 的次序读取标准表，把数字位于第一行记为一步向上，位于第二行记为一步向下，路径从高度零出发。列严格递增等价于每个前缀的向上步数不少于向下步数。总路径数为 \(\binom{a+b}{b}\)。对第一次越到高度 \(-1\) 的坏路径，把到该点为止的上、下步互换，就得到从高度零出发、具有 \(a+1\) 个上步和 \(b-1\) 个下步的路径，其终点高度为 \(a-b+2\ge2\)。反过来，每条这样的路径必经过高度 \(1\)；把到第一次达到高度 \(1\) 为止的上、下步互换，便得到第一次在该位置越到高度 \(-1\) 的坏路径。这两次反射互逆，因此
\[
f^{(a,b)}=\binom{a+b}{b}-\binom{a+b}{b-1}.
\]
当 \(a=b=r\) 时，这是 Catalan 数
\[
f^{(r,r)}=\frac1{r+1}\binom{2r}{r}.
\]
该推导与\cref{b5:thm:hook-length}给出的两行钩长计算一致，但它还解释了相减的那项在数什么。
